\documentclass[11pt]{report} %size of font and type of layout
\usepackage[a4paper,margin=1.1in]{geometry}

\usepackage{lscape}
\usepackage{pgfplots}
\usepackage{tikz}
\usepackage{xfrac}
\usepackage[super]{nth} %allow 1st, 2nd by typing \nth{1}, \nth{2}
\usepackage[title,titletoc,page]{appendix} %for using appendices
\usepackage{graphicx} %have images
\usepackage{chngcntr} %be able to make figures go 1, 2, 3 instead of 1.1, 1.2, 1.3, etc
\usepackage{subcaption} %have subfigures
\usepackage{wrapfig} %be able to wrap text around figures
\usepackage{amsmath} %a way to typeset maths
\usepackage{amssymb} %able to type maths symbols
\usepackage{textcomp} %so we can use \textdegree to get degree symbol
\usepackage[parfill]{parskip} %remove paragraph indentation
\usepackage{setspace} %be able to set line spacing
\usepackage{tocloft} %more control over table of contents, list of figures, etc
\usepackage{ltxtable} %get tabularx split over pages
\usepackage[english]{babel} %english-specific charaters and language rules
\usepackage{url} %type urls
\usepackage{parskip} %space between paragraphs
\usepackage[UKenglish]{isodate} %change date format
%\usepackage[hang,flushmargin]{footmisc} %no indents in footers
%\usepackage{verbatimbox} %able to use addvbuffer to change spaces before and after tables
\usepackage{pdfpages} %insert PDF pages
%\usepackage{titletoc} %create a small table of contents at the start of each section
\usepackage{microtype} %improves spacing between words and letters
\usepackage[nodayofweek]{datetime} %use \today to get the date
\usepackage[inline]{enumitem} %have inline lists
\usepackage{fancyhdr}		%for getting customized headers
\usepackage{lastpage}		%enable getting the total of the pages
%\usepackage{floatrow} %use for indicating figures' sources
%\usepackage{showframe} %show margins on page
\usepackage[hidelinks]{hyperref} %hyperlinks in pdf
\usepackage{cleveref} %automatic referencing
\usepackage{bookmark}
\DeclareGraphicsExtensions{.pdf,.png,.jpg,.eps} %can just give images' names
\usepackage{epstopdf} %allows for encapsulated postscript vector images, better quality and scalable
\usepackage[none]{hyphenat}		%kills all word breaks that were in annoying places

%--------------------------heading spacings----------------------------%
\usepackage{titlesec} %be able to change title settings
	\titleformat{\chapter}{\normalfont\huge\bfseries}{\thechapter.}{1em}{} %have number before chapter title
	\titlespacing{\chapter}{0pc}{0pc}{0pc} %have nice spacing after chapter title {indent} {before} {after}
	\titlespacing{\section}{0pc}{1pc}{0pc} %have nice spacing after section title {indent} {before} {after}
	\titlespacing{\subsection}{1pc}{0pc}{0pc} %have nice spacing after subsection title {indent} {before} {after}
	\titlespacing{\subsubsection}{0pc}{0pc}{-1pc} %have nice spacing after subsection title {indent} {before} {after}
	\titlelabel{\thetitle.\quad} %full stop after the number in the heading title
\usepackage{etoolbox} %remove page break before chapter heading
	\makeatletter
	\patchcmd{\chapter}{\thispagestyle{plain}}{\thispagestyle{fancy}}{}{} %allows the header/footer on chapter pages
	\patchcmd{\chapter}{\if@openright\cleardoublepage\else\clearpage\fi}{}{}{}
	\makeatother
%----------------------------figures & lof (list of figures)-----------------------------%
\renewcommand{\cftfigfont}{Figure } %list of figures includes "Figure"
\newcommand*{\noaddvspace}{\renewcommand*{\addvspace}[1]{}}\addtocontents{lof}{\protect\noaddvspace} %equal line spacing in the list of figures
\setlength{\cftfigindent}{0pt} %remove indentation from figures in lof
\usepackage{float} %able to use H to force figures to go here
%\floatstyle{boxed}	%this will place a thin border around my figures
%\restylefloat{figure}
\AtBeginEnvironment{figure}{\vspace{10pt}} %spacing before figure
\AtEndEnvironment{figure}{\vspace{-10pt}} %spacing after figure
\setlength{\abovecaptionskip}{10pt} %spacing before figure caption
\setlength{\belowcaptionskip}{5pt} %spacing after figure caption
\usepackage[justification=centering,font={normalsize, it}]{caption} %control caption text
%----------------------------tables & lot------------------------------%
\renewcommand{\cfttabfont}{Table } %list of tables includes "Table"
\setlength{\cfttabindent}{0pt} %remove indentation from tables in lot
\usepackage{tabularx} %word wrap in tables
\usepackage{longtable} %allow tables to go over multiple pages
\usepackage{multirow} %make cells span multiple rows in tables
\captionsetup{belowskip=5pt,aboveskip=4pt}
\newcolumntype{L}{>{\centering\arraybackslash}m{3cm}} %centre and wrap text in tables
\usepackage{booktabs} %use \toprule, \midrule, \bottomrule and \cmidrule for better spacing in tables
\usepackage{pbox} %force line break in a table cell
%---------------------------create own lists---------------------------%
\newlistof{example}{exp}{List of Examples} %create own lists
\newcommand{\example}[1]{ %create own lists
	\refstepcounter{example} %create own lists
	\par\noindent\textbf{Example \theexample. #1} %create own lists
	\addcontentsline{exp}{example} %create own lists
	{\protect\numberline{\thechapter.\theexample}#1}\par} %create own lists
%---------------------------------font---------------------------------%
%\sfdefault = calibri-like font, \rmdefault = times new roman-like
\renewcommand\UrlFont{\rmfamily}
\renewcommand{\familydefault}{\rmdefault}
%----------------------------------------------------------------------%
\begin{document}
\begin{onehalfspace} %single line spacing
%\counterwithout{figure}{chapter} %makes figures go 1, 2, 3 instead of 1.1, 1.2, 1.3, etc
%\counterwithout{table}{chapter} %makes tables go 1, 2, 3 instead of 1.1, 1.2, 1.3, etc
%\counterwithout{section}{chapter} %make sections 1, 2, 3 not 1.1 etc
\pgfplotsset{compat=1.12} %because it complained
\pretolerance=10000 %disable hyphenating at line breaks
\abovedisplayskip=5pt %remove space before equations
\belowdisplayskip=5pt %remove space after equations
\cleanlookdateon %remove current date format to replace with [UKenglish]{isodate}
%\let\savenumberline\numberline\def\numberline#1{\savenumberline{#1.}} %adds dot after chapter title in toc, comment this out to get appendices to work and not cause spacing errors.
%\AtEndEnvironment{itemize}{\vspace{-5pt}} %spacing after lists 

\newcommand{\thetitle}	{Business Plan}
\newcommand{\theauthor}	{Roberto Aldera - ALDROB001 \\ Gareth Callanan - CLLGAR010 \\ Luke Goemans - GMNLUK001 \\ Benjamin Scholtz - SCHBEN011 \\ Munawwar Tayob - TYBMUN001 }
\newcommand{\studentnum}{}
\newcommand{\thedate}	{\today}
\newcommand{\coursecode}{EEE4051F}
\newcommand{\coursename}{New Venture Planning}
\newcommand{\staffmember}{Dr. Peter Martinez}

\pagenumbering{gobble}	%no page number for title page
\renewcommand\headrulewidth{0pt}	%get rid of stray header underline
	\begin{center}
\begin{minipage}{1\textwidth}
\begin{center}
	
\begin{center}
\includegraphics[width=1\linewidth]{"UCThoriz"}
\end{center}
\vskip20pt
\Large{Faculty of Engineering and the Built Environment}\\
Department of Electrical Engineering\\
\coursecode
\vskip10pt
\begin{figure}[H]
\centering
\hspace{-3em}
\includegraphics[width=0.6\textwidth]{images/uvuka_logo_ben}
\vskip20pt
\end{figure}
\hrule
\vskip20pt
\Huge\sc{\thetitle}
\vskip20pt
\hrule
\vskip20pt
\Large\textnormal{\textbf{Team 5:}\\
\theauthor\\
\vskip10pt
\textbf{Lecturer:} \staffmember
\vskip20pt
\hrule
\vskip20pt
\thedate\\}
%\textbf{\assignmentname}
\end{center}
\end{minipage}
\end{center}

\newpage 	
%	\chapter*{Plagiarism Declaration}
We know that plagiarism is wrong. Plagiarism is to use another's work and pretend that it is our own.

We have used the IEEE convention for citation and referencing. In this report, all contributions to, and quotations from, the work(s) of other people have been cited and referenced. 

This report is our own work. We have not allowed, and will not allow, anyone to copy our work.
\vskip100pt
\begin{center}
	Signed: \line(1,0){200}
	\vskip20pt
	Dated: \line(1,0){200}
\end{center}

\newpage
	\input{table_of_contents.tex}
\renewcommand\headrulewidth{0.4pt}
\pagenumbering{arabic}	%start the page numbering (and reset it to zero)

\pagestyle{fancy}
\fancyhf{}
\rhead{Group 5}
\lhead{\coursecode \\ \coursename}
\cfoot{Page \thepage \ of \pageref{LastPage}}

%\chapter*{Executive Summary}
%\addcontentsline{toc}{chapter}{Executive Summary}

\newpage
\chapter{Introduction}
Ladies and gentlemen! We want your money. This is a reference example\cite{nise2011control}. Add more text here\cite{ExampleForWebsiteReference}.

\newpage
\chapter{Product Description}
\section{Product Overview}
Uvuka seeks to reduce road accidents through intelligent monitoring of vehicle drivers. The product consists of a single device that mounts onto the dashboard of a vehicle, as seen in \cref{fig:Uvuka_casing_concept}. The basic working version of the product will have a camera pointing in the direction of the driver that focuses on the driver's eyes. The device will monitor a driver's eyes and detect if they close. If the eyes close for a set period of time (longer than the average blink), then the device assumes that the driver is sleeping and takes corrective measures to wake the driver up.

\begin{figure}[H]
\centering
\includegraphics[width=0.6\textwidth]{images/Uvuka_casing_concept.JPG}
\vskip10pt
\caption[The caption that appears in the list of figures]{Conceptual design of Uvuka device rendered in SolidWorks}
\label{fig:Uvuka_casing_concept}
\end{figure}

Uvuka consists of a range of products. This product range will all center around the same goal (keeping the driver awake) but with different features present in different models. The first prototype produced will be the most elementary device, and once orders have been secured and units delivered,  more complex products can be developed and sold.

\section{Technical description of Uvuka}
\label{sec:Technical description of Uvuka}
The basic product consists of four main hardware components:
\begin{enumerate}
  \item \textit{Cabin Camera} - The cabin camera faces the driver and provides a real-time video feed of their face, upon which algorithms can be performed. The camera is able to capture suitable low light footage for all driving conditions.
  \item \textit{Audio Speaker} - The audio speaker is used to alert the driver. The speaker will play randomly generated tones that vary in frequency and duration. This prevents the driver from growing accustomed to the tones and thus will be more likely to wake him up\cite{Habituation}.
  \item \textit{Micro-controller and Inertial Navigation System} - The micro-controller needs to be powerful enough in order to perform real-time video processing. Eye-tracking and detection is done using OpenCV algorithms - these are free and open source video and image processing software algorithms. The micro-controller will also control the speaker and begin to play sounds when the driver begins to fall asleep \cite{OpenCV}. An Inertial Navigation System (INS) will include the necessary accelerometers, gyroscopes, and magnetometers to record the vehicle's dynamics and driver's handling of the vehicle.
  \item \textit{Power Supply} - The device can either be plugged into the standard 12 V cigarette plug, or it can be directly connected to the car power source for permanent mounts. The power is regulated to the correct voltage within the device, which will also include surge protection.
\end{enumerate}

The product has various modes of operation as seen in \cref{fig:DeviceFlowChart} below. The system defaults to the passive scanning mode where the device detects the driver's face and monitors eye movement. If the eye is determined to be closed, the device  enters a minor emergency state and assumes the driver's concentration has begun to wane. A brief sound is emitted to alert the driver. If this does not cause a satisfactory level of alertness, it assumes the driver is at risk of falling asleep and proceeds to take more drastic action. The action taken depends on the time the driver remains unresponsive and will vary from one type of device to the other.

\begin{figure}[H]
\centering
\includegraphics[width=0.4\textwidth]{images/ProductFlowchart}
\vskip10pt
\caption[Device Operation Activity Diagram]{Activity Diagram of the device operation}
\label{fig:DeviceFlowChart}
\end{figure}

\section{Product Range}
The the device prototype explained in \cref{sec:Technical description of Uvuka} above is envisioned as being the entry level product that private consumers will be able to purchase. This will be targeted for those on a tighter budget who require this device. As the business grows, the product range will be expanded based on consumer feedback and market research.

There are a number of other products and expansions planned:
\begin{itemize}
\item Frequent Driver Uvuka camera - This product is aimed at taxi operators and delivery drivers that are behind the wheel for extended periods of time. They are more likely to drive while tired or near exhausted (some taxi drivers can have shifts over 14 hours long \cite{taxiDriverHours}). In this case, sound may not be enough to wake them up. This product will have a strip attached to the steering wheel that will vibrate if the driver begins to fall asleep.
\item Long Distance Transport Uvuka camera - this product also incorporates the cabin camera and vibrating steering wheel attachment. Additionally, it also has a second camera facing the road in front of the driver and a GPS unit to track the vehicle's location. All this data is recorded on an SD card. This system is aimed for long distance freight vehicles and allows fleet managers to monitor their drivers and vehicles through the INS data for quality and insurance purposes.
\end{itemize}

\section{Competitor product analysis}
There are currently two primary international competitors producing a similar product. Seeing Machines, a company based in Australia, have developed a technology that tracks driver eyelid movement through computer vision algorithms\cite{SeeingMachinesWebsite}. They monitor the frequency and duration of blinking, as well as eyelid velocity. If the driver’s head begins to drop, an audio alarm is triggered and the seat vibrates. If they fall asleep again, a dispatcher is alerted and may make radio contact with the driver to check if they need a break. Seeing Machines has specifically worked with drivers that operate Caterpillar mining equipment\cite{SeeingMachinesWired}.

The second competitor company is Exeros, based in the United Kingdom\cite{Exeros}. Their product is targeted at businesses that operate a fleet of vehicles. When the driver begins to fall asleep, audio tones are sounded to attempt to wake them up. The devices they produce are able to be connected to an external GPS so that when an alarm is triggered, an outside individual can be notified of the occurrence and the location.

% --Extra sources/info--
% Toyota has implemented a driver alertness sensor as early as 2008 in their upmarket Toyota Crown cars. This camera detects the upper and lower eyelids, calculating how open they are. The car has built-in proximity sensors outside the car to measure what’s happening and to apply the brakes or prepare the airbags depending on the probability of a collision.

% HARMAN technology has developed a pupil-based driver monitoring system which interprets the level of the driver’s cognitive load and mental multitasking ability, adapting the car’s interface and connected devices accordingly. This drowsiness monitoring can cause the car to automatically turn the driver’s mobile device into do-not-disturb mode or adjusting ADAS system intervention thresholds to minimize physical and mental distraction to the driver.
% IQ Intel

\section{Advantages of Uvuka over competitor products} 
Uvuka’s unique features allow it to perform several of the tasks currently performed by competing products.
%TODO expand on this section^^^

\section{Product positioning}
The product will be versatile enough to cater for the needs of a number of industries. The following sectors are considered:
\begin{itemize}
\item Heavy duty freight truck drivers
\item Bus drivers
\item Construction equipment operators (cranes, diggers, etc.)
\item Taxi drivers
\item Private use
\end{itemize}

\section{Intellectual Property and freedom of trade}
%Luke wants to do this section :P

\newpage
\chapter{Marketing Plan}
\section{Problem background}
Over 5500 people lost their lives in 4500 accidents on South Africa's roads between 2014 and 2015. This harrowing statistic stirred in us a desire to engineer a product capable of making our roads a safer place for all\cite{EWNRoadDeaths}.

South Africa's railway system is currently unable to support the needs of the growing economy. With 88\% of all freight being moved on the roads\cite{BDlive_freight}, many truck drivers are exposed daily to the risks of falling asleep at the wheel\cite{ArriveAliveDriverTiredness}.

\section{Uvuka product logo}
In order to convey the product's purpose and the company's values, a logo was designed. This seeks to motivate the product to potential investors and customers as a device that keeps an eye on you in order to ensure your safety on the roads.

\begin{figure}[H]
\centering
\includegraphics[width=0.7\textwidth]{images/uvuka_logo_ben}
\vskip10pt
\caption[The caption that appears in the list of figures]{The Uvuka logo}
\label{fig:uvuka_logo}
\end{figure}

\section{Market research}
No local competitors currently exist in the market. We expect to capture a significant portion of the market due to the product’s relative novelty and necessity.

\begin{appendices}
\appendix
\chapter{Insert Here}
\label{app:Insert Here}	
\newpage

\chapter{Another}
\label{app:Another}
\end{appendices}	%for use with appendices
\newpage
\renewcommand\bibname{References}
\bibliography{Aldera_bibliography}
\bibliographystyle{ieeetr}
%changing this to ^ieeetran^ and downloading the zip and compiling on Rob's PC will produce a reference list that has the URLs in it, don't worry about this for now.
\addcontentsline{toc}{chapter}{References}

\newpage	%this keeps the header/footer working on the last page. Do not type anything after this.

\end{onehalfspace}
\end{document}

%-------------------------------Figures--------------------------------%
\begin{figure}[H]
\centering
\includegraphics[width=0.4\textwidth]{images/%}
\vskip10pt
\caption[The caption that appears in the list of figures]{The caption that appears under the figure}
\label{fig:exam}
\end{figure}

\begin{figure}[H]
\centering
\begin{subfigure}[t]{0.49\textwidth}
\centering
\includegraphics[width=0.9\linewidth]{images/%}
\caption{This is figure a}
\label{fig:figureA}
\end{subfigure}
\begin{subfigure}[t]{0.49\textwidth}
\centering
\includegraphics[width=0.9\linewidth]{images/%}
\caption{This is figure b}
\label{fig:figureB}
\end{subfigure}
\caption{Caption}
\end{figure}

\begin{wrapfigure}{r}{0.5\textwidth}
\begin{center}
\includegraphics[width=0.48\textwidth]{images/ %}
\end{center}
\caption{A gull}
\end{wrapfigure}

As can be seen in \cref{fig:exam}\\
%-------------------------------Headings-------------------------------%
\chapter{Chapter Title}
\label{ch:Chapter Title}

\section{Section Title}
\label{sec:Section Title}

\subsection{Subsection Title}
\label{subsec:Subsection Title}

\subsubsection{Subsubsection Title}
\label{bb:Subsubsection Title}
%---------------------------Tables and Lists---------------------------%
\input{\filepath /table_of_contents.tex}
\input{\filepath /list_of_figures.tex}
\input{\filepath /list_of_tables.tex}
\input{\filepath /list_of_example.tex}

\startcontents[chapters] %mini table of contents
\printcontents[chapters]{}{1}{}
\newpage
%-----------------------------Insert Pages-----------------------------%
\includepdf[pages={1}]{ %filename.pdf}
%-------------------------------Appendix-------------------------------%
\startcontents[section]
\printcontents[section]{}{1}{}
\newpage
\addcontentsline{toc}{section}{1. %}
%--------------------------------Citing--------------------------------%
\footnote{\cite{aa}}
\cite{aa,bb}
%--------------------------Custom Bibliography-------------------------%
makebst
%----------------------------Maths & Symbols---------------------------%
http://www.artofproblemsolving.com/Wiki/index.php/LaTeX:Symbols
http://www.tex.ac.uk/tex-archive/info/symbols/comprehensive/symbols-a4.pdf
$Hardcore \: maths \: here \: \therefore \ddot{\theta}^{2} = \left(\dfrac{\omega_{n}}{17a}\right) \times 9$
$$ More \: maths: \sqrt{52} + F_{1} = 6 \times 10^{5} \: Hz$$
$\mathrm{\today}$ %maths in normal font
%---------------------------Lines and Spacing--------------------------%
\noindent\makebox[\linewidth]{\rule{\textwidth}{0.4pt}}
\vspace{11pt}
\hspace{-11pt}
\hphantom{11pt}
%----------------------------Page Numbering----------------------------%
\pagenumbering{gobble} %remove page numbering
\pagenumbering{roman} %roman letters page numbering
\pagenumbering{arabic} %arabic letters page numbering
\setcounter{page}{1} %reset page numbers
%--------------------------------Tables--------------------------------%
\begin{table}[H]
\caption{Caption of table}
\addvbuffer[12pt 8pt]{
\begin{tabularx}{0.3\textwidth}{@{}rcl@{}}
\hline
\multicolumn{3}{@{}c@{}}{\textbf{Title}}\\
\hline
Info	& 3	& mm\\
Item	& 5	& x $\approx$ 34\\
\cline{2-2}
Width	& 4	& miles\\
Depth	& 7	& \pbox{20cm}{This is the first \\ cell}\\
\hline
\label{tab:capt}
\end{tabularx}
}
\end{table}

\begin{tabularx}{\textwidth}{@{}lX@{}}
 & \\
\end{tabularx}

\\
\\
\begin{tabular}{ l l }
 & \\
\end{tabular}
%--------------------------------Lists---------------------------------%
http://en.wikibooks.org/wiki/LaTeX/List_Structures

%TC:ignore
\begin{itemize} \itemsep1pt \parskip0pt \parsep0pt \vspace{-6pt}
%TC:endignore
\item[-] The first item
\item[-] The second item
\item[-] The third etc \ldots
\end{itemize}

\begin{enumerate} [label=\bfseries Exercise \arabic*:]
\item The first item
\item The second item
\item The third etc \ldots
\end{enumerate}

\begin{description}
\item[First] The first item
\item[Second] The second item
\item[Third] The third etc \ldots
\end{description}

\begin{description}
\item[First] \hfill \\
The first item
\item[Second] \hfill \\
The second item
\item[Third] \hfill \\
The third etc \ldots
\end{description}

Inline lists are \begin{enumerate*}[label=\itshape\alph*\upshape)]
\item formatted within their paragraph;
\item usually labelled with letters; and
\item usually have the final item prefixed with
`and' or `or',
\end{enumerate*} like this example.

\example{Your first example}
\label{1st_ex}
\example{Your second example}
\label{2nd_ex}
\example{Your third example. (See example \ref{1st_ex} and \ref{2nd_ex})}
%----------------------------------------------------------------------%
\fi